\begin{corollary}[黄金分支在好侧未达单箱最优置中且保有显式正缺口]\label{cor:xi-golden-w1-good-side-gap-to-single-box-optimum}\leanverified{paper\_xi\_golden\_w1\_good\_side\_gap\_to\_single\_box\_optimum}
沿用定理 \ref{thm:xi-kronecker-w1-single-box-optimal-centering} 与
定理 \ref{thm:xi-golden-w1-true-two-phase-limit} 的记号。取
\[
\alpha=\varphi^{-1},
\qquad
q=F_n,
\qquad
p=F_{n-1},
\qquad
\delta_n:=\alpha-\frac{F_{n-1}}{F_n}
=
\frac{(-1)^{n-1}}{\sqrt5\,F_n^2}\Bigl(1-(-1)^n\varphi^{-2n}\Bigr).
\]
记 \(\mathsf{T}_{F_n}^{\min}\) 为定理 \ref{thm:xi-kronecker-w1-single-box-optimal-centering} 在 \(q=F_n\) 处给出的理论最小值。则在好侧子序列 \(n\) 为奇数时：
\begin{enumerate}
  \item 存在 \(n_0\) 使得对一切奇数 \(n\ge n_0\)，都有
  \[
  \delta_n<\delta_\ast(F_n).
  \]
  因而黄金分支的真实箱内偏移始终落在好侧抛物线的下降支上，尚未达到单箱最优置中点。
  \item 存在显式正极限缺口
  \[
  \boxed{\;
  \lim_{\substack{n\to\infty\\ n\ \mathrm{odd}}}
  F_n\Bigl(\mathsf{T}_{F_n}(\varphi^{-1})-\mathsf{T}_{F_n}^{\min}\Bigr)
  =
  \frac{61}{240}-\frac{1}{2\sqrt5}
  >0.\;}
  \]
\end{enumerate}
\end{corollary}
\begin{proof}
当 \(n\) 为奇数时，由引理 \ref{lem:golden-convergent-error-closed}，
\[
F_n^2\delta_n
=
\frac{1}{\sqrt5}\bigl(1+\varphi^{-2n}\bigr)
\longrightarrow \frac1{\sqrt5}.
\]
另一方面，
\[
\delta_\ast(F_n)=\frac{3}{2F_n(2F_n-1)},
\qquad
F_n^2\delta_\ast(F_n)=\frac{3F_n}{2(2F_n-1)}
\longrightarrow \frac34.
\]
由于 \(1/\sqrt5<3/4\)，故最终必有 \(\delta_n<\delta_\ast(F_n)\)，得到第一条。

再由定理 \ref{thm:xi-kronecker-w1-single-box-optimal-centering},
\[
F_n\,\mathsf{T}_{F_n}^{\min}
=
\frac12-\frac{3(F_n-1)}{8(2F_n-1)}
\longrightarrow
\frac12-\frac{3}{16}
=
\frac{5}{16}.
\]
而定理 \ref{thm:xi-golden-w1-true-two-phase-limit} 给出
\[
F_n\,\mathsf{T}_{F_n}(\varphi^{-1})
\longrightarrow
\frac12-\frac{1}{2\sqrt5}+\frac1{15}
\qquad
(n\to\infty,\ n\ \mathrm{odd}).
\]
两式相减便得
\[
\lim_{\substack{n\to\infty\\ n\ \mathrm{odd}}}
F_n\Bigl(\mathsf{T}_{F_n}(\varphi^{-1})-\mathsf{T}_{F_n}^{\min}\Bigr)
=
\left(\frac12-\frac{1}{2\sqrt5}+\frac1{15}\right)-\frac{5}{16}
=
\frac{61}{240}-\frac{1}{2\sqrt5}.
\]
该常数严格为正，故得到第二条。证毕。
\end{proof}

\endinput
